% Luc Destombe --- licence CC BY-NC-SA 4.0
% https://creativecommons.org/licenses/by-nc-sa/4.0/deed.fr
% Fichier autonome : compiler avec pdflatex, deux fois (signets, nombre de pages).
\documentclass[11pt,a4paper]{article}
\makeatletter
\newif\iflycee@unecolonne
\lycee@unecolonnetrue
\newif\iflycee@palatino
\lycee@palatinotrue

\RequirePackage[utf8]{inputenc}
\RequirePackage[T1]{fontenc}
\RequirePackage[french]{babel}
\RequirePackage{etoolbox}

\newcommand{\ShorthandsOff}{\@ifundefined{shorthandoff}{}{\shorthandoff{;:!?}}}
\newcommand{\ShorthandsOn}{\@ifundefined{shorthandon}{}{\shorthandon{;:!?}}}
\ShorthandsOff
\AtBeginDocument{\ShorthandsOn}
\AtBeginEnvironment{tikzpicture}{\ShorthandsOff}

\RequirePackage{geometry}
\RequirePackage{fancyhdr}
\RequirePackage{lastpage}
\RequirePackage{multicol}
\RequirePackage{array}
\RequirePackage{enumitem}
\RequirePackage{xcolor}
\RequirePackage{needspace}

\RequirePackage{amsmath,amssymb}
\RequirePackage{mathpazo}
\DeclareMathSizes{10.95}{10.95}{8}{5.5}
\begingroup\catcode`\_=\active \catcode`\^=\active
\gdef_#1{\sb{\scriptscriptstyle #1}}
\gdef^#1{\sp{\scriptscriptstyle\lycee@moinsepais #1}}
\endgroup
\newcommand\lycee@moins{\mathbin{%
  \setbox\z@\hbox{$\m@th\scriptscriptstyle\lycee@moinsorig$}%
  \dimen@=.55\wd\z@ \advance\dimen@-.5pt
  \hbox to.75\wd\z@{\hss$\m@th\scriptscriptstyle\vcenter{\hbox{%
    \tikz\draw[line width=.5pt,line cap=round](0,0)--(\the\dimen@,0);}}$\hss}}}
\begingroup\lccode`\~=`\- \lowercase{\endgroup\let~}\lycee@moins
\newcommand\lycee@moinsepais{\mathcode`\-="8000 }
\AtBeginDocument{\mathchardef\lycee@moinsorig=\mathcode`\-\relax}
\AtBeginDocument{\catcode`\_=12 \mathcode`\_="8000
  \catcode`\^=12 \mathcode`\^="8000 }
\newcommand\lycee@exposants{%
  \fontdimen13\textfont2=.48\fontdimen6\textfont2
  \fontdimen14\textfont2=.48\fontdimen6\textfont2
  \fontdimen15\textfont2=.48\fontdimen6\textfont2
  \fontdimen13\scriptfont2=.48\fontdimen6\scriptfont2
  \fontdimen14\scriptfont2=.48\fontdimen6\scriptfont2
  \fontdimen15\scriptfont2=.48\fontdimen6\scriptfont2}
\every@math@size\expandafter{\the\every@math@size\lycee@exposants}
\newlength{\retraitcalculs}
\setlength{\retraitcalculs}{2em}

\RequirePackage{tikz}
\usetikzlibrary{arrows.meta,calc,babel,shapes.misc}
\RequirePackage[most]{tcolorbox}
\tcbuselibrary{skins,breakable}

\definecolor{coulDef}{HTML}{1F4E79}
\definecolor{coulProp}{HTML}{7030A0}
\definecolor{coulRem}{HTML}{C55A11}
\definecolor{coulEx}{HTML}{2E7D32}
\definecolor{coulExo}{HTML}{B03A2E}
\definecolor{coulPart}{HTML}{1F4E79}
\definecolor{coulMeth}{HTML}{1B6E4F}
\definecolor{coulCourbe}{HTML}{0B5ED7}
\definecolor{coulCourbeB}{HTML}{C00000}
\definecolor{coulCourbeC}{HTML}{2E7D32}
\definecolor{coulGrille}{HTML}{8C8C8C}
\definecolor{coulRep}{HTML}{1B6E4F}
\definecolor{coulBar}{HTML}{2E7D32}

\newcommand{\R}{\mathbb{R}}
\newcommand{\Rs}{\mathbb{R}^{*}}
\renewcommand{\le}{\leqslant}
\renewcommand{\ge}{\geqslant}
\renewcommand{\approx}{\simeq}
\newcommand{\vect}[1]{\overrightarrow{#1}}
\newcommand{\norme}[1]{\left\lVert #1 \right\rVert}
\newcommand{\intervalle}[2]{\left[#1\,;#2\right]}
\RequirePackage{cancel}
\renewcommand{\CancelColor}{\color{coulExo}}
\newcommand{\fois}[1]{\times{\color{coulExo}#1}}
\newcommand{\barre}[1]{\cancel{#1}}

\tikzset{grille/.style={coulGrille,line width=0.4pt}}
\newcommand{\quadrillage}[4]{\draw[grille,step=1] (#1,#2) grid (#3,#4);}
\tikzset{
  croix/.style={cross out,draw,line width=0.6pt,inner sep=0pt,minimum size=4pt},
  croixdroite/.style={croix,rotate=45,line width=0.8pt},
}
\newcommand{\pt}[3]{\path (#1) node[croix]{} node[#2,font=\small,inner sep=2.5pt]{$#3$};}

\newcommand{\lycee@repere}[4]{%
  \draw[grille,step=1] (#1,#2) grid (#3,#4);
  \draw[-{Stealth[length=2mm]},thick] (#1,0)--({#3+0.4},0) node[below left=-1pt]{$x$};
  \draw[-{Stealth[length=2mm]},thick] (0,#2)--(0,{#4+0.4}) node[below left=-1pt]{$y$};
  \node[below left,font=\scriptsize,inner sep=1.5pt] at (0,0) {$0$};}
\newcommand{\lycee@gradx}[1]{\foreach \k/\l in {#1}{%
  \draw (\k,0.07)--(\k,-0.07) node[below,font=\scriptsize,inner sep=1.5pt]{$\l$};}}
\newcommand{\lycee@grady}[1]{\foreach \k/\l in {#1}{%
  \draw (0.07,\k)--(-0.07,\k) node[left,font=\scriptsize,inner sep=1.5pt]{$\l$};}}
\newcommand{\lycee@intff}[2]{\left[#1\,;#2\right]}
\newcommand{\lycee@intfo}[2]{\left[#1\,;#2\right[}
\newcommand{\lycee@intof}[2]{\left]#1\,;#2\right]}
\newcommand{\lycee@intoo}[2]{\left]#1\,;#2\right[}
\newcommand{\lycee@ens}[1]{\left\{#1\right\}}
\AtBeginDocument{%
  \providecommand{\repere}{\lycee@repere}%
  \providecommand{\gradx}{\lycee@gradx}%
  \providecommand{\grady}{\lycee@grady}%
  \providecommand{\Cf}{\mathcal{C}\sb{\scriptscriptstyle f}}%
  \providecommand{\Cg}{\mathcal{C}\sb{\scriptscriptstyle g}}%
  \providecommand{\intff}{\lycee@intff}%
  \providecommand{\intfo}{\lycee@intfo}%
  \providecommand{\intof}{\lycee@intof}%
  \providecommand{\intoo}{\lycee@intoo}%
  \providecommand{\ens}{\lycee@ens}}

\newlist{questions}{enumerate}{3}
\setlist[questions,1]{label=\arabic*\textdegree),leftmargin=*,itemsep=2pt,topsep=3pt}
\setlist[questions,2]{label=\alph*),leftmargin=*,itemsep=1pt,topsep=2pt}
\setlist[questions,3]{label=\roman*),leftmargin=*,itemsep=1pt,topsep=1pt}
\setlist[itemize]{label=\textbullet,leftmargin=*,itemsep=2pt,topsep=3pt}

\newcommand{\besoinplace}[1]{\par\penalty\@M\begingroup
  \dimen@=#1\relax\dimen@ii=\pagegoal\advance\dimen@ii-\pagetotal
  \ifdim\dimen@>\dimen@ii\ifdim\dimen@ii>\z@\vfil\fi\break\fi\endgroup}

\newcommand{\lycee@pied}{\thepage/\pageref{LastPage}}
\newcommand{\entete}[2]{%
  \pagestyle{fancy}\fancyhf{}%
  \fancyhead[L]{\footnotesize\color{coulPart}\textbf{#1}}%
  \fancyhead[R]{\footnotesize\color{coulPart}\textbf{#2}}%
  \fancyfoot[C]{\footnotesize\lycee@pied}%
  \renewcommand{\headrulewidth}{0.4pt}%
  \renewcommand{\headrule}{\hbox to\headwidth{%
    \color{coulPart!40}\leaders\hrule height \headrulewidth\hfill}}}

\RequirePackage{ccicons}
\newcommand{\lycee@licence}{{\scriptsize\color{gray}%
  \copyright\ Luc Destombe --- \ccbyncsa\ CC BY-NC-SA 4.0}}
\newif\iflycee@licence \lycee@licencetrue
\newcommand{\sanslicence}{\lycee@licencefalse}
\AtBeginDocument{\iflycee@licence\AddToHookNext{shipout/foreground}{%
  \put(0,0){\rlap{\kern\dimexpr1in+\hoffset+\oddsidemargin+\textwidth\relax
    \llap{\raisebox{-\dimexpr1in+\voffset+\topmargin+\headheight+\headsep
      +\textheight+\footskip\relax}{\lycee@licence}}}}}\fi}

\newcommand{\formulesserrees}{%
  \setlength{\abovedisplayskip}{3pt}\setlength{\belowdisplayskip}{3pt}%
  \setlength{\abovedisplayshortskip}{2pt}\setlength{\belowdisplayshortskip}{3pt}}

\setlength{\parindent}{0pt}

\RequirePackage[implicit=false,hidelinks,pdfencoding=auto,psdextra,bookmarksopen,
  bookmarksopenlevel=1,pdfcreator={},pdfproducer={}]{hyperref}
\RequirePackage{bookmark}
\hypersetup{pdfauthor={Luc Destombe},pdfkeywords={CC BY-NC-SA 4.0}}
\newcounter{lycee@signet}
\newcommand{\signet}[3][]{\stepcounter{lycee@signet}%
  \edef\lycee@dest{\if\relax\detokenize{#1}\relax
    signet.\arabic{lycee@signet}\else#1\fi}%
  \hypertarget{\lycee@dest}{}\bookmark[level=#2,dest=\lycee@dest]{#3}}
\pdfstringdefDisableCommands{%
  \def\R{R}\def\vect#1{#1}\def\Cf{Cf}\def\Cg{Cg}\def\fois#1{×#1}%
  \def\barre#1{#1}\def\intervalle#1#2{[#1;#2]}\def\intff#1#2{[#1;#2]}%
  \def\textsuperscript#1{#1}\def\dfrac#1#2{#1/#2}\def\frac#1#2{#1/#2}%
  \def\vec#1{#1⃗}}%
\RequirePackage[official]{eurosym}

\tcbset{exercice corrige/.style={
  enhanced,breakable,before upper=\raggedright,
  colback=white,colframe=coulExo,
  boxrule=0.8pt,arc=2pt,
  left=6pt,right=6pt,top=8pt,bottom=5pt,
  before skip=9pt,after skip=4pt,
  coltitle=white,fonttitle=\bfseries\small,
  attach boxed title to top left={xshift=8pt,yshift=-\tcboxedtitleheight/2},
  boxed title style={colback=coulExo,arc=2pt,boxrule=0pt}}}
\newcounter{exo}
\newtcolorbox{exercice}[1][]{exercice corrige,
  phantom={\signet[exercice-\arabic{exo}]{2}{Exercice \arabic{exo}}},
  title={Exercice \theexo\ifx\relax#1\relax\else{} --- #1\fi}}
\newenvironment{exo}[1][\relax]{\refstepcounter{exo}\begin{exercice}[#1]}{\end{exercice}}

\geometry{top=1.8cm,bottom=1cm,left=1cm,right=1cm,headheight=14pt,footskip=0.6cm}
\renewcommand{\lycee@pied}{\thepage}

\renewcommand{\theexo}{\ifnum\value{exo}<10 0\fi\arabic{exo}}
\tcbset{exercice corrige/.style={
  enhanced,breakable,before upper=\raggedright,
  colback=white,colframe=coulExo,
  boxrule=0.9pt,arc=2pt,
  left=8pt,right=8pt,top=8pt,bottom=6pt,
  before skip=8pt,after skip=6pt,
  coltitle=white,fonttitle=\bfseries,
  attach boxed title to top left={xshift=8pt,yshift=-\tcboxedtitleheight/2},
  boxed title style={colback=coulExo,arc=2pt,boxrule=0pt}}}

\newcommand{\entetefiche}[2]{%
  \begin{center}
    \begin{tcolorbox}[enhanced,width=0.92\linewidth,colback=white,colframe=coulPart,
        boxrule=1.4pt,arc=3pt,halign=center,top=6pt,bottom=6pt]
      {\large\bfseries\color{coulPart} CORRIGÉ --- FICHE D'EXERCICES}\\[3pt]
      {\normalsize\bfseries #1}\\[2pt]
      {\normalsize #2}
    \end{tcolorbox}
  \end{center}}

\newcounter{partie}
\newcommand{\partie}[1]{%
  \stepcounter{partie}%
  \besoinplace{8\baselineskip}\vspace{4pt}%
  \begin{tcolorbox}[enhanced,colback=coulPart!8,colframe=coulPart,boxrule=0pt,
      leftrule=3pt,arc=0pt,left=6pt,right=4pt,top=3pt,bottom=3pt,
      before skip=6pt,after skip=2pt]
    \bfseries\color{coulPart}\Roman{partie} --- #1\signet{1}{\Roman{partie} --- #1}
  \end{tcolorbox}}
\newcommand{\partiesansnum}[1]{%
  \besoinplace{8\baselineskip}\vspace{4pt}%
  \begin{tcolorbox}[enhanced,colback=coulPart!8,colframe=coulPart,boxrule=0pt,
      leftrule=3pt,arc=0pt,left=6pt,right=4pt,top=3pt,bottom=3pt,
      before skip=6pt,after skip=2pt]
    \bfseries\color{coulPart}#1\signet{1}{#1}
  \end{tcolorbox}}

\setlist[questions,1]{label=\arabic*\textdegree),leftmargin=*,itemsep=3pt,topsep=3pt}
\setlist[questions,2]{label=\alph*),leftmargin=*,itemsep=2pt,topsep=2pt}
\setlist[questions,3]{label=\roman*),leftmargin=*,itemsep=1pt,topsep=1pt}
\newlist{expr}{enumerate}{1}
\setlist[expr]{label={\Alph*\ $=$},leftmargin=*,itemsep=2pt,topsep=3pt}

\newcommand{\res}[1]{{\color{coulRep}#1}}
\newcommand{\rem}[1]{\par\smallskip{\small\itshape\color{black!60}$\rightarrow$ #1}}
\newenvironment{calculs}{\par\smallskip\noindent$\displaystyle\begin{aligned}[t]}%
  {\end{aligned}$\par\smallskip}
\newcommand{\justif}[1]{\quad\text{\small\itshape\color{black!60}#1}}

\setlength{\parskip}{2pt}

\makeatother
\geometry{top=1.8cm,bottom=1.6cm,left=1.8cm,right=1.8cm,headheight=14pt,footskip=30pt}
\entete{Chapitre 5 --- Suites numériques --- corrigé}{Première spécialité}
\usepackage{tkz-tab}

\newtcblisting{python}{%
  listing only,enhanced,
  colback=black!3,colframe=black!35,boxrule=0.5pt,arc=1pt,
  left=4pt,right=4pt,top=1pt,bottom=1pt,
  listing options={language=Python,basicstyle=\ttfamily\small,
    keywordstyle=\color{coulMeth}\bfseries,columns=fullflexible,
    keepspaces=true,showstringspaces=false,aboveskip=0pt,belowskip=0pt}}
\newcommand{\reperesuite}[4]{%
  \draw[grille,step=1] (-1,#2) grid (#1,#3);
  \draw[-{Stealth[length=2mm]},thick] (-1,0)--({#1+0.4},0)
    node[below left=-1pt,font=\footnotesize]{$n$};
  \draw[-{Stealth[length=2mm]},thick] (0,#2)--(0,{#3+0.4})
    node[below left=-1pt,font=\footnotesize]{$#4$};
  \node[below left,font=\scriptsize,inner sep=1.5pt] at (0,0) {$0$};}

\begin{document}

\entetefiche{Chapitre 5 --- Suites numériques}{Première spécialité mathématiques}

\partie{Définition d'une suite}

\begin{exo}[Repérer les termes et les rangs]
On numérote à partir de $0$ : $u_{0}=5$ ; $u_{1}=9$ ; $u_{2}=13$ ; $u_{3}=17$ ;
$u_{4}=21$ ; $u_{5}=25$.
\begin{questions}[label=\arabic*.]
  \item $\res{u_{0}=5}$ et $\res{u_{3}=17}$.
  \item $21=u_{4}$ : le terme $21$ est de $\res{\text{rang } 4}$.
  \item Le 4\textsuperscript{e} terme est $\res{u_{3}=17}$ (on compte $u_{0}$, $u_{1}$,
        $u_{2}$, $u_{3}$).
  \item Le terme suivant est $\res{u_{4}=21}$ ; le précédent est $\res{u_{2}=13}$.
\end{questions}
\end{exo}

\begin{exo}[Une suite qui commence à $u_{1}$]
\begin{questions}[label=\arabic*.]
  \item $\res{v_{1}=1}$ et $\res{v_{4}=16}$.
  \item Le premier terme est $v_{1}$ : le 5\textsuperscript{e} terme est
        $\res{v_{5}=25}$, de rang $5$.
  \item Ce sont les carrés des rangs : on conjecture $\res{v_{n}=n^{2}}$.
\end{questions}
\end{exo}

\begin{exo}[Premier rang]
\begin{questions}[label=\alph*)]
  \item Il faut $n\neq 0$ : $\res{n_{0}=1}$.
  \item Il faut $n-3\ge 0$, soit $n\ge 3$ : $\res{n_{0}=3}$.
  \item Il faut $n-5>0$ (racine et division), soit $n>5$ : $\res{n_{0}=6}$.
  \item Il faut $n\ge 0$ (racine) et $n\neq 0$ (division) : $\res{n_{0}=1}$.
\end{questions}
\end{exo}

\begin{exo}[Les décimales de $\frac{1}{7}$]
\begin{questions}[label=\arabic*.]
  \item $u_{3}=2$ ; $u_{4}=8$ ; $u_{5}=5$ ; $u_{6}=7$ ; $u_{7}=1$ ; $u_{8}=4$. Le premier
        rang est $\res{1}$ ($u_{1}$ est la première décimale).
  \item Les chiffres $1$ ; $4$ ; $2$ ; $8$ ; $5$ ; $7$ se répètent tous les $6$ rangs.
        $20=3\times 6+2$, donc $\res{u_{20}=u_{2}=4}$ ; $100=16\times 6+4$, donc
        $\res{u_{100}=u_{4}=8}$.
\end{questions}
\rem{Le reste de la division par $6$ donne la position dans le motif.}
\end{exo}

\partie{Suites définies explicitement}

\begin{exo}[Calculer des termes]
\begin{questions}[label=\alph*)]
  \item $u_{0}=-3$ ; $u_{1}=1$ ; $u_{5}=4\times 5-3=\res{17}$.
  \item $v_{0}=3$ ; $v_{1}=2$ ; $v_{5}=25-10+3=\res{18}$.
  \item $w_{0}=\frac{2}{1}=2$ ; $w_{1}=\frac{5}{2}$ ; $w_{5}=\res{\frac{17}{6}}$.
  \item $t_{0}=2^{0}-1=0$ ; $t_{1}=1$ ; $t_{5}=32-1=\res{31}$.
\end{questions}
\end{exo}

\begin{exo}[Retrouver un rang]
\begin{questions}[label=\arabic*.]
  \item $u_{10}=3\times 10+7=\res{37}$.
  \item \begin{calculs}
3n+7 &= 52 &&\justif{on résout}\\
3n &= 45\\
n &= 15
\end{calculs}
        $\res{u_{15}=52}$.
  \item $3n+7=101$ donne $3n=94$, soit $n=\frac{94}{3}$, qui n'est pas entier :
        $\res{101 \text{ n'est pas un terme}}$.
\end{questions}
\end{exo}

\begin{exo}[Exprimer $u_{n+1}$ et $u_{n-1}$]
\begin{questions}[label=\alph*)]
  \item $u_{n+1}=5(n+1)-2=\res{5n+3}$ ; $u_{n-1}=5(n-1)-2=\res{5n-7}$.
  \item \begin{calculs}
v_{n+1} &= (n+1)^{2}-4(n+1) &&\justif{on remplace $n$ par $(n+1)$}\\
v_{n+1} &= n^{2}+2n+1-4n-4\\
v_{n+1} &= n^{2}-2n-3\\[3pt]
v_{n-1} &= (n-1)^{2}-4(n-1)\\
v_{n-1} &= n^{2}-2n+1-4n+4\\
v_{n-1} &= n^{2}-6n+5
\end{calculs}
        $\res{v_{n+1}=n^{2}-2n-3}$ et $\res{v_{n-1}=n^{2}-6n+5}$.
  \item $w_{n+1}=\dfrac{(n+1)+3}{2(n+1)+1}=\res{\dfrac{n+4}{2n+3}}$ ;
        $w_{n-1}=\dfrac{(n-1)+3}{2(n-1)+1}=\res{\dfrac{n+2}{2n-1}}$.
\end{questions}
\end{exo}

\begin{exo}[Ne pas confondre]
\begin{questions}[label=\alph*)]
  \item \begin{calculs}
u_{n+1} &= (n+1)^{2}+2(n+1)\\
u_{n+1} &= n^{2}+2n+1+2n+2 = \res{n^{2}+4n+3}
\end{calculs}
  \item $u_{n}+1=\res{n^{2}+2n+1}$.
  \item $u_{2n}=(2n)^{2}+2\times 2n=\res{4n^{2}+4n}$.
  \item $2u_{n}=\res{2n^{2}+4n}$.
\end{questions}
\rem{Quatre expressions différentes : on remplace $n$ (a, c) ou l'on opère sur le
terme (b, d).}
\end{exo}

\partie{Suites définies par récurrence}

\begin{exo}[Calculer de proche en proche]
\begin{questions}[label=\alph*)]
  \item $u_{1}=2\times 3-1=5$ ; $u_{2}=2\times 5-1=9$ ; $u_{3}=2\times 9-1=\res{17}$.
  \item $v_{1}=v_{0}+3\times 0-1=1$ ; $v_{2}=v_{1}+3\times 1-1=3$ ;
        $v_{3}=v_{2}+3\times 2-1=\res{8}$.
  \item $w_{2}=\frac{16}{2}+4=12$ ; $w_{3}=\frac{12}{2}+4=10$ ;
        $w_{4}=\frac{10}{2}+4=\res{9}$.
\end{questions}
\rem{En b), le $n$ de la relation est le rang du terme précédent.}
\end{exo}

\begin{exo}[Avec un carré, avec une fraction]
\begin{questions}[label=\alph*)]
  \item $a_{1}=1^{2}+1=2$ ; $a_{2}=2^{2}+1=5$ ; $a_{3}=5^{2}+1=\res{26}$.
  \item $b_{1}=\frac{1}{1+1}=\frac{1}{2}$ ; $b_{2}=\frac{1}{1+\frac{1}{2}}=\frac{2}{3}$ ;
        $b_{3}=\frac{1}{1+\frac{2}{3}}=\res{\frac{3}{5}}$.
\end{questions}
\end{exo}

\begin{exo}[Traduire une situation]
\begin{questions}[label=\alph*)]
  \item $u_{0}=1\,800$ et $\res{u_{n+1}=1{,}02\,u_{n}}$ ($+2\,\%$ : on multiplie par
        $1{,}02$).
  \item $u_{0}=12\,000$ et $\res{u_{n+1}=0{,}95\,u_{n}+900}$ (il reste $95\,\%$, puis on
        ajoute les arrivants).
  \item $u_{0}=500$ et $\res{u_{n+1}=u_{n}-40}$.
\end{questions}
\end{exo}

\begin{exo}[Une forêt]
\begin{questions}[label=\arabic*.]
  \item $u_{1}=5\,000-0{,}08\times 5\,000+300=5\,000-400+300=\res{4\,900}$ ;
        $u_{2}=0{,}92\times 4\,900+300=4\,508+300=\res{4\,808}$.
  \item Abattre $8\,\%$, c'est multiplier par $1-0{,}08=0{,}92$ ; on ajoute ensuite les
        $300$ arbres plantés : $\res{u_{n+1}=0{,}92\,u_{n}+300}$.
\end{questions}
\end{exo}

\partie{Calculatrice}

\begin{exo}[Afficher des termes]
\begin{questions}[label=\arabic*.]
  \item \texttt{Explicite}, \texttt{Tableau} de $0$ à $10$ :
\begin{center}
\renewcommand{\arraystretch}{1.2}\small
\begin{tabular}{|c|*{11}{c|}}
\hline
$n$ & $0$ & $1$ & $2$ & $3$ & $4$ & $5$ & $6$ & $7$ & $8$ & $9$ & $10$\\\hline
$u_{n}$ & $30$ & $21$ & $14$ & $9$ & $6$ & $5$ & $6$ & $9$ & $14$ & $21$ & $30$\\\hline
\end{tabular}
\end{center}
  \item Par exemple $X$ de $0$ à $11$ et $Y$ de $0$ à $32$.
  \item \texttt{Récurrente d'ordre 1} : $\res{v_{10}=1\,026}$ et
        $\res{v_{20}=1\,048\,578}$.
\end{questions}
\end{exo}

\begin{exo}[Conjecturer une expression]
\begin{questions}[label=\arabic*.]
  \item $w_{0}$ à $w_{6}$ : $1$ ; $2$ ; $5$ ; $10$ ; $17$ ; $26$ ; $37$. Puis
        $\res{w_{10}=101}$ et $\res{w_{50}=2\,501}$.
  \item Chaque terme est un carré plus $1$ : $\res{w_{n}=n^{2}+1}$.
\end{questions}
\end{exo}

\begin{exo}[Représenter et conjecturer]
\begin{questions}[label=\arabic*.]
  \item $20$ ; $10$ ; $6{,}67$ ; $5$ ; $4$ ; $3{,}33$ ; $2{,}86$ ; $2{,}5$ ; $2{,}22$ ; $2$.
  \item Les points descendent : la suite semble $\res{\text{décroissante}}$.
\end{questions}
\end{exo}

\begin{exo}[Le carré ou la puissance ?]
\begin{questions}[label=\arabic*.]
  \item
\begin{center}
\renewcommand{\arraystretch}{1.2}\small
\begin{tabular}{|c|*{11}{c|}}
\hline
$n$ & $0$ & $1$ & $2$ & $3$ & $4$ & $5$ & $6$ & $7$ & $8$ & $9$ & $10$\\\hline
$n^{2}$ & $0$ & $1$ & $4$ & $9$ & $16$ & $25$ & $36$ & $49$ & $64$ & $81$ & $100$\\\hline
$2^{n}$ & $1$ & $2$ & $4$ & $8$ & $16$ & $32$ & $64$ & $128$ & $256$ & $512$ & $1\,024$\\\hline
\end{tabular}
\end{center}
  \item $u_{n}>v_{n}$ seulement pour $\res{n=3}$ ; $u_{n}=v_{n}$ pour $\res{n=2}$ et
        $\res{n=4}$.
  \item On conjecture que, pour $n\ge 5$, $\res{2^{n}>n^{2}}$ : l'écart grandit vite.
\end{questions}
\end{exo}

\partie{Programmer en Python}

\begin{exo}[Écrire la fonction]
\begin{minipage}[t]{0.47\linewidth}
\begin{python}
def terme_v(n):
    v = 4
    for i in range(n):
        v = 3*v - 5
    return v
\end{python}
\end{minipage}\hfill
\begin{minipage}[t]{0.47\linewidth}
\texttt{terme\_v(10)} renvoie $\res{88\,576}$.
\end{minipage}
\end{exo}

\begin{exo}[Lire un programme]
\begin{questions}[label=\arabic*.]
  \item La suite définie par $\res{u_{0}=5}$ et $\res{u_{n+1}=\frac{u_{n}}{2}+3}$ : la
        fonction renvoie $u_{n}$.
  \item $\texttt{mystere(1)}=\frac{5}{2}+3=\res{5{,}5}$ ;
        $\texttt{mystere(2)}=\frac{5{,}5}{2}+3=\res{5{,}75}$ ;
        $\texttt{mystere(3)}=\frac{5{,}75}{2}+3=\res{5{,}875}$.
\end{questions}
\end{exo}

\begin{exo}[Une relation qui dépend de $n$]
\begin{questions}[label=\arabic*.]
  \item $w_{1}=0+0+3=3$ ; $w_{2}=3+2+3=8$ ; $w_{3}=8+4+3=\res{15}$.
  \item Ligne à compléter : \texttt{w = w + 2*i + 3}.
  \item \texttt{terme\_w(10)} renvoie $\res{120}$.
\end{questions}
\rem{À la première boucle, \texttt{i} vaut $0$ : on calcule bien $w_{1}$ à partir de
$w_{0}$.}
\end{exo}

\begin{exo}[La suite de Fibonacci]
\begin{questions}[label=\arabic*.]
  \item $f_{2}=2$ ; $f_{3}=3$ ; $f_{4}=5$ ; $f_{5}=8$ ; $\res{f_{6}=13}$.
  \item Ligne à compléter : \texttt{a, b = b, a + b} (le nouveau terme est la somme des
        deux précédents).
  \item \texttt{fibo(20)} renvoie $\res{10\,946}$.
\end{questions}
\end{exo}

\partie{Représentation graphique}

\begin{exo}[Représenter une suite récurrente]
\begin{minipage}[c]{0.58\linewidth}\raggedright
\begin{questions}[label=\arabic*.]
  \item $u_{0}=2$ ; $u_{1}=3$ ; $u_{2}=3{,}5$ ; $u_{3}=3{,}75$ ; $u_{4}=3{,}875$ ;
        $u_{5}=3{,}9375$.
  \item Voir le graphique.
  \item Les points montent : la suite semble $\res{\text{croissante}}$ (et se rapprocher
        de $4$).
\end{questions}
\end{minipage}\hfill
\begin{minipage}[c]{0.38\linewidth}\centering
\begin{tikzpicture}[x=0.6cm,y=0.6cm]
  \reperesuite{6}{0}{5}{u_{n}}
  \gradx{1,2,3,4,5} \grady{1,2,3,4}
  \foreach \n/\u in {0/2,1/3,2/3.5,3/3.75,4/3.875,5/3.9375}
    \path[coulCourbe] (\n,\u) node[croix]{};
\end{tikzpicture}
\end{minipage}
\end{exo}

\begin{exo}[Représenter une suite explicite]
\begin{minipage}[c]{0.58\linewidth}\raggedright
\begin{questions}[label=\arabic*.]
  \item $4{,}5$ ; $2$ ; $0{,}5$ ; $0$ ; $0{,}5$ ; $2$ ; $4{,}5$.
  \item Voir le graphique.
  \item La suite $\res{\text{décroît jusqu'au rang } 3}$, puis
        $\res{\text{croît}}$.
\end{questions}
\end{minipage}\hfill
\begin{minipage}[c]{0.38\linewidth}\centering
\begin{tikzpicture}[x=0.5cm,y=0.5cm]
  \reperesuite{7}{0}{5}{v_{n}}
  \gradx{1,2,3,4,5,6} \grady{1,2,3,4}
  \foreach \n in {0,1,...,6}
    \path[coulCourbe] (\n,{(\n-3)^2/2}) node[croix]{};
\end{tikzpicture}
\end{minipage}
\end{exo}

\begin{exo}[Lire un nuage de points]
\begin{questions}[label=\arabic*.]
  \item $\res{u_{0}=1}$ ; $\res{u_{2}=4}$ ; $\res{u_{5}=1}$.
  \item $u_{n}=4$ pour $\res{n=2}$ et $\res{n=3}$.
  \item La suite croît de $u_{0}$ à $u_{2}$, reste égale à $4$ de $u_{2}$ à $u_{3}$, puis
        décroît jusqu'à $u_{6}=0$.
\end{questions}
\end{exo}

\begin{exo}[Des points sur deux droites]
\begin{minipage}[c]{0.58\linewidth}\raggedright
\begin{questions}[label=\arabic*.]
  \item $0$ ; $-1$ ; $2$ ; $-3$ ; $4$ ; $-5$ ; $6$.
  \item Ni l'un ni l'autre : $u_{1}<u_{0}$ mais $u_{2}>u_{1}$.
  \item Rangs pairs : $u_{n}=n$, points sur $\res{y=x}$ ; rangs impairs : $u_{n}=-n$,
        points sur $\res{y=-x}$.
\end{questions}
\end{minipage}\hfill
\begin{minipage}[c]{0.38\linewidth}\centering
\begin{tikzpicture}[x=0.32cm,y=0.32cm]
  \reperesuite{7}{-6}{7}{u_{n}}
  \gradx{2,4,6} \grady{-4,-2,2,4,6}
  \draw[coulCourbe!50,dashed] (0,0)--(6.8,6.8) (0,0)--(6.8,-6.8);
  \foreach \n in {0,1,...,6}
    \path[coulCourbe] (\n,{(-1)^\n*\n}) node[croix]{};
\end{tikzpicture}
\end{minipage}
\end{exo}

\partie{Sens de variation : signe de $u_{n+1}-u_{n}$}

\begin{exo}[Étudier le signe de la différence]
\begin{questions}[label=\alph*)]
  \item \begin{calculs}
u_{n+1}-u_{n} &= 3(n+1)^{2}+(n+1)-(3n^{2}+n)\\
u_{n+1}-u_{n} &= 3n^{2}+6n+3+n+1-3n^{2}-n\\
u_{n+1}-u_{n} &= 6n+4
\end{calculs}
        $6n+4>0$ : $(u_{n})$ est $\res{\text{strictement croissante}}$.
  \item $v_{n+1}-v_{n}=n^{2}+2>0$ : $(v_{n})$ est $\res{\text{strictement croissante}}$.
  \item \begin{calculs}
w_{n+1}-w_{n} &= (n+1)^{2}-10(n+1)-(n^{2}-10n)\\
w_{n+1}-w_{n} &= n^{2}+2n+1-10n-10-n^{2}+10n\\
w_{n+1}-w_{n} &= 2n-9
\end{calculs}
        $2n-9\ge 0$ pour $n\ge 4{,}5$ : $(w_{n})$ est $\res{\text{strictement croissante
        à partir du rang } 5}$.
\end{questions}
\end{exo}

\begin{exo}[Suites décroissantes]
\begin{questions}[label=\alph*)]
  \item $u_{n+1}-u_{n}=7-2(n+1)^{2}-7+2n^{2}=-4n-2<0$ : $\res{\text{strictement
        décroissante}}$.
  \item $v_{n+1}-v_{n}=-(n-3)^{2}\le 0$ : $\res{\text{décroissante}}$ (avec
        $v_{4}=v_{3}$, elle n'est pas strictement décroissante).
  \item \begin{calculs}
w_{n+1}-w_{n} &= \frac{1}{n+2}-\frac{1}{n+1}\\
w_{n+1}-w_{n} &= \frac{1\fois{(n+1)}}{(n+2)\fois{(n+1)}}-\frac{1\fois{(n+2)}}{(n+1)\fois{(n+2)}} &&\justif{même dénominateur}\\
w_{n+1}-w_{n} &= \frac{-1}{(n+1)(n+2)}
\end{calculs}
        Un numérateur négatif et un dénominateur positif : $\res{\text{strictement
        décroissante}}$.
\end{questions}
\end{exo}

\begin{exo}[Conjecturer, puis démontrer]
\begin{questions}[label=\arabic*.]
  \item $u_{1}=13$ ; $u_{2}=8$ ; $u_{3}=5$ ; $u_{4}=4$ ; $u_{5}=5$ : la suite semble
        décroître, puis croître à partir du rang $4$.
  \item $u_{n+1}-u_{n}=2n-7$, positif pour $n\ge 3{,}5$.
  \item La suite est $\res{\text{croissante à partir du rang } 4}$.
\end{questions}
\end{exo}

\begin{exo}[Un carré toujours positif]
\begin{questions}[label=\arabic*.]
  \item $v_{n+1}-v_{n}=(v_{n}-1)^{2}\ge 0$ (un carré) : $\res{(v_{n}) \text{ est
        croissante}}$.
  \item $v_{1}=0+1=1$ ; $v_{2}=1+0=1$ ; $v_{3}=1$. La suite est constante à partir du
        rang $1$ : $\res{\text{elle n'est pas strictement croissante}}$.
\end{questions}
\end{exo}

\partie{Sens de variation : quotient $\dfrac{u_{n+1}}{u_{n}}$}

\begin{exo}[Comparer le quotient à 1]
Tous les termes sont strictement positifs.
\begin{questions}[label=\alph*)]
  \item $\dfrac{u_{n+1}}{u_{n}}=\dfrac{2\times 5^{n+1}}{2\times 5^{n}}=5>1$ :
        $\res{\text{strictement croissante}}$.
  \item $\dfrac{v_{n+1}}{v_{n}}=0{,}8<1$ : $\res{\text{strictement décroissante}}$.
  \item \begin{calculs}
\frac{w_{n+1}}{w_{n}} &= \frac{4^{n+1}}{n+1}\times\frac{n}{4^{n}} &&\justif{multiplier par l'inverse}\\
\frac{w_{n+1}}{w_{n}} &= \frac{4n}{n+1}\\
\frac{w_{n+1}}{w_{n}}-1 &= \frac{4n-n-1}{n+1}=\frac{3n-1}{n+1}
\end{calculs}
        Pour $n\ge 1$ on a $3n-1>0$ : $\res{\text{strictement croissante}}$.
\end{questions}
\end{exo}

\begin{exo}[D'autres quotients]
\begin{questions}[label=\alph*)]
  \item $\dfrac{u_{n+1}}{u_{n}}=\dfrac{3}{2}=1{,}5>1$ : $\res{\text{strictement
        croissante}}$.
  \item $\dfrac{v_{n+1}}{v_{n}}=\dfrac{n+1}{2^{n+1}}\times\dfrac{2^{n}}{n}=\dfrac{n+1}{2n}$
        et $\dfrac{n+1}{2n}-1=\dfrac{1-n}{2n}\le 0$ : $\res{\text{décroissante}}$
        ($v_{1}=v_{2}=\frac{1}{2}$).
  \item $\dfrac{w_{n+1}}{w_{n}}=\dfrac{5(n+1)}{n+2}$ et
        $\dfrac{5(n+1)}{n+2}-1=\dfrac{4n+3}{n+2}>0$ : $\res{\text{strictement
        croissante}}$.
\end{questions}
\end{exo}

\begin{exo}[Choisir la méthode]
\begin{questions}[label=\alph*)]
  \item Quotient : $0{,}5<1$ et termes positifs : $\res{\text{strictement
        décroissante}}$.
  \item Différence : $v_{n+1}-v_{n}=2n+6>0$ : $\res{\text{strictement croissante}}$.
  \item Quotient : $\dfrac{(n+1)\times 2^{n+1}}{n\times 2^{n}}=\dfrac{2(n+1)}{n}>1$ :
        $\res{\text{strictement croissante}}$.
\end{questions}
\end{exo}

\begin{exo}[Croissante à partir d'un rang]
\begin{questions}[label=\arabic*.]
  \item $u_{1}=2$ ; $u_{2}=1$ ; $u_{3}=\frac{8}{9}\approx 0{,}89$ ; $u_{4}=1$ ;
        $u_{5}=1{,}28$.
  \item $\dfrac{u_{n+1}}{u_{n}}=\dfrac{2^{n+1}}{(n+1)^{2}}\times\dfrac{n^{2}}{2^{n}}
        =\dfrac{2n^{2}}{(n+1)^{2}}$. Ce quotient est $\ge 1$ si
        $2n^{2}\ge (n+1)^{2}$, soit $2n^{2}\ge n^{2}+2n+1$, soit $n^{2}-2n-1\ge 0$.
  \item $\Delta=8$ ; racines $1-\sqrt{2}$ et $1+\sqrt{2}\approx 2{,}41$. Pour $n$ entier
        $\ge 1$, $n^{2}-2n-1\ge 0$ dès $n\ge 3$ : $\res{\text{croissante à partir du
        rang } 3}$.
\end{questions}
\end{exo}

\partie{Sens de variation : fonction associée}

\begin{exo}[Étudier la fonction associée]
\begin{questions}[label=\alph*)]
  \item $f(x)=x^{3}+6x^{2}+1$ ; $f'(x)=3x^{2}+12x=3x(x+4)\ge 0$ sur
        $\intfo{0}{+\infty}$ : $\res{(u_{n}) \text{ strictement croissante}}$.
  \item $f(x)=\dfrac{2x+3}{x+1}$ ;
        $f'(x)=\dfrac{2(x+1)-(2x+3)}{(x+1)^{2}}=\dfrac{-1}{(x+1)^{2}}<0$ :
        $\res{(v_{n}) \text{ strictement décroissante}}$.
  \item $f(x)=x^{3}-3x^{2}$ ; $f'(x)=3x^{2}-6x=3x(x-2)$ ; $f(2)=-4$.
\begin{center}
\begin{tikzpicture}[scale=0.8]
  \tkzTabInit[lgt=2,espcl=2.4]{$x$/0.8,$f'(x)$/0.8,$f(x)$/1.4}{$0$,$2$,$+\infty$}
  \tkzTabLine{z,-,z,+,}
  \tkzTabVar{+/$0$,-/$-4$,+/{}}
\end{tikzpicture}
\end{center}
        $\res{(w_{n}) \text{ croissante à partir du rang } 2}$.
\end{questions}
\end{exo}

\begin{exo}[D'autres fonctions]
\begin{questions}[label=\alph*)]
  \item $f'(x)=2x-8$ : $f$ décroît sur $\intff{0}{4}$, croît sur
        $\intfo{4}{+\infty}$ : $\res{(u_{n}) \text{ croissante à partir du rang } 4}$.
  \item $f(x)=5-\dfrac{4}{x+2}$ ; $f'(x)=\dfrac{4}{(x+2)^{2}}>0$ :
        $\res{(v_{n}) \text{ strictement croissante}}$.
  \item $f'(x)=\dfrac{(x^{2}+1)-x\times 2x}{(x^{2}+1)^{2}}=\dfrac{1-x^{2}}{(x^{2}+1)^{2}}$,
        négatif pour $x\ge 1$ : $\res{(w_{n}) \text{ décroissante à partir du rang }
        1}$.
\end{questions}
\end{exo}

\begin{exo}[Une erreur fréquente]
\begin{questions}[label=\arabic*.]
  \item $u_{1}=0{,}5\times 10+2=7$ ; $u_{2}=5{,}5$ ; $u_{3}=4{,}75$.
  \item La suite décroît : $\res{\text{l'affirmation est fausse}}$. La propriété ne vaut
        que pour une suite explicite $u_{n}=f(n)$.
\end{questions}
\end{exo}

\begin{exo}[Suite croissante, fonction non croissante]
\begin{questions}[label=\arabic*.]
  \item $f'(x)=2(x-0{,}4)$ : $f$ décroît sur $\intff{0}{0{,}4}$ et croît sur
        $\intfo{0{,}4}{+\infty}$.
  \item \begin{calculs}
u_{n+1}-u_{n} &= (n+0{,}6)^{2}-(n-0{,}4)^{2}\\
u_{n+1}-u_{n} &= (n+0{,}6-n+0{,}4)(n+0{,}6+n-0{,}4) &&\justif{$a^{2}-b^{2}=(a-b)(a+b)$}\\
u_{n+1}-u_{n} &= 2n+0{,}2
\end{calculs}
        $2n+0{,}2>0$ : $\res{(u_{n}) \text{ strictement croissante}}$.
  \item $\res{\text{Non}}$ : ici la suite croît alors que $f$ n'est pas croissante sur
        $\intfo{0}{+\infty}$. La propriété du cours ne se renverse pas.
\end{questions}
\end{exo}

\partie{Notion de limite}

\begin{exo}[Conjecturer une limite]
\begin{questions}[label=\alph*)]
  \item $u_{9}=2{,}5$ ; $u_{99}=2{,}05$ ; $u_{999}=2{,}005$ : limite $\res{2}$.
  \item $v_{10}=301$ ; $v_{100}=30\,001$ ; $v_{1\,000}=3\,000\,001$ : limite
        $\res{+\infty}$.
  \item $4$ ; $-4$ ; $4$ ; $-4$ : $\res{\text{pas de limite}}$.
\end{questions}
\end{exo}

\begin{exo}[Limite finie ou infinie]
\begin{questions}[label=\alph*)]
  \item $-95$ ; $-9\,995$ ; $-999\,995$ : limite $\res{-\infty}$.
  \item $0{,}01$ ; $0{,}000\,1$ ; $0{,}000\,001$ : limite $\res{0}$.
  \item $2{,}1$ ; $2{,}01$ ; $2{,}001$ : limite $\res{2}$.
\end{questions}
\end{exo}

\begin{exo}[Avec la calculatrice]
\begin{questions}[label=\arabic*.]
  \item $u_{10}\approx 5{,}995\,117$ ; $u_{20}\approx 5{,}999\,995$ ;
        $u_{30}\approx 6{,}000\,000$.
  \item On conjecture $\res{\lim\limits_{n\to+\infty}u_{n}=6}$.
\end{questions}
\end{exo}

\begin{exo}[Alterner et converger]
\begin{questions}[label=\arabic*.]
  \item $u_{10}=0{,}1$ ; $u_{11}\approx -0{,}091$ ; $u_{100}=0{,}01$ ;
        $u_{101}\approx -0{,}009\,9$.
  \item Le signe change à chaque rang, mais les termes se rapprochent de $0$ : on
        conjecture $\res{\lim\limits_{n\to+\infty}u_{n}=0}$.
\end{questions}
\rem{Une suite qui alterne peut avoir une limite : il suffit que les écarts à $0$
deviennent aussi petits que l'on veut.}
\end{exo}

\partie{Problèmes : déterminer un seuil}

\begin{exo}[Abonnés d'une newsletter]
\begin{questions}[label=\arabic*.]
  \item $u_{1}=0{,}85\times 3\,000+800=\res{3\,350}$ ;
        $u_{2}=0{,}85\times 3\,350+800=\res{3\,647{,}5}$.
  \item $15\,\%$ partent : il reste $0{,}85\,u_{n}$ ; on ajoute $800$.
  \item \begin{minipage}[t]{0.45\linewidth}\vspace{0pt}
\begin{python}
def seuil():
    u = 3000
    n = 0
    while u <= 5000 :
        u = 0.85*u + 800
        n = n + 1
    return n
\end{python}
\end{minipage}
  \item $u_{11}\approx 4\,942{,}9$ et $u_{12}\approx 5\,001{,}4$ : $\res{\text{au bout de }
        12 \text{ mois}}$.
\end{questions}
\end{exo}

\begin{exo}[Seuil d'une suite explicite]
\begin{questions}[label=\arabic*.]
  \item \texttt{while u <= 1000 :} puis \texttt{u = 2*n**2 + 3*n} (on calcule
        directement le terme de rang \texttt{n}).
  \item $u_{21}=945$ et $u_{22}=1\,034$ : $\res{n=22}$.
\end{questions}
\end{exo}

\begin{exo}[Élimination d'un médicament]
\begin{questions}[label=\arabic*.]
  \item $m_{1}=32$ ; $m_{2}=25{,}6$ ; $\res{m_{n+1}=0{,}8\,m_{n}}$.
  \item \begin{minipage}[t]{0.45\linewidth}\vspace{0pt}
\begin{python}
def seuil():
    m = 40
    n = 0
    while m >= 5 :
        m = 0.8*m
        n = n + 1
    return n
\end{python}
\end{minipage}
  \item $m_{9}\approx 5{,}37$ et $m_{10}\approx 4{,}29$ : $\res{10 \text{ heures}}$.
\end{questions}
\end{exo}

\begin{exo}[Qui dépasse l'autre ?]
\begin{questions}[label=\arabic*.]
  \item Tant que $b_{n}\le a_{n}=500+40n$, on passe au mois suivant : la fonction
        renvoie le premier mois où le second site a plus d'abonnés que le premier.
  \item $b_{19}\approx 1\,223<a_{19}=1\,260$ et $b_{20}\approx 1\,345>a_{20}=1\,300$ :
        elle renvoie $\res{20}$ ; au bout de $20$ mois, le site à $+10\,\%$ dépasse
        l'autre.
\end{questions}
\end{exo}

\partiesansnum{Exercices type bac}

\begin{exo}[Automatismes (QCM)]
$\res{1.\ b)}$ : $(n+1)^{2}-1=n^{2}+2n$ ;\quad $\res{2.\ b)}$ : $u_{1}=5$ puis
$u_{2}=14$ ;\quad $\res{3.\ c)}$ ;\quad $\res{4.\ b)}$ : la différence est
négative ;\quad $\res{5.\ b)}$.
\end{exo}

\begin{exo}[Une plateforme de covoiturage]
\begin{questions}[label=\arabic*.]
  \item $u_{1}=0{,}9\times 600+40=\res{580}$ ; $u_{2}=0{,}9\times 580+40=\res{562}$. Il
        reste $90\,\%$ des conducteurs, puis $40$ s'ajoutent : $u_{n+1}=0{,}9\,u_{n}+40$.
  \item $n=0$ : $400+200=600=u_{0}$ ; $n=1$ : $400+180=580=u_{1}$.
  \item \begin{calculs}
u_{n+1}-u_{n} &= 200\times 0{,}9^{n+1}-200\times 0{,}9^{n} &&\justif{les $400$ s'éliminent}\\
u_{n+1}-u_{n} &= 200\times 0{,}9^{n}(0{,}9-1) &&\justif{facteur commun}\\
u_{n+1}-u_{n} &= -20\times 0{,}9^{n}
\end{calculs}
        $0{,}9^{n}>0$, donc la différence est négative : $\res{(u_{n}) \text{ est
        strictement décroissante}}$.
  \item $0{,}9^{n}$ se rapproche de $0$ : $\res{\lim\limits_{n\to+\infty}u_{n}=400}$. Le
        nombre de conducteurs se stabiliserait vers $400$.
  \item \texttt{while u >= 450 :} et \texttt{u = 0.9*u + 40}. $u_{13}\approx 450{,}8$ et
        $u_{14}\approx 445{,}8$ : la fonction renvoie $14$, soit $\res{\text{mars 2027}}$.
\end{questions}
\end{exo}

\begin{exo}[Coût moyen de fabrication]
\begin{questions}[label=\arabic*.]
  \item $c_{1}=\res{401}$ ; $c_{10}=10+40=\res{50}$ ; $c_{20}=20+20=\res{40}$.
  \item \begin{calculs}
f'(x) &= 1-\frac{400}{x^{2}} &&\justif{$\left(\frac{1}{x}\right)'=-\frac{1}{x^{2}}$}\\
f'(x) &= \frac{x^{2}-400}{x^{2}}=\frac{(x-20)(x+20)}{x^{2}}
\end{calculs}
\begin{center}
\begin{tikzpicture}[scale=0.8]
  \tkzTabInit[lgt=2,espcl=2.4]{$x$/0.8,$f'(x)$/0.8,$f(x)$/1.4}{$1$,$20$,$+\infty$}
  \tkzTabLine{,-,z,+,}
  \tkzTabVar{+/$401$,-/$40$,+/{}}
\end{tikzpicture}
\end{center}
  \item $(c_{n})$ décroît jusqu'au rang $20$, puis croît : le coût moyen est le plus
        faible pour $\res{20 \text{ lampes}}$ ($40$~€).
  \item Pour $n>0$, $c_{n}>50$ revient à $n^{2}+400>50n$, soit $n^{2}-50n+400>0$.
        $\Delta=900$ ; racines $10$ et $40$ ; le trinôme est positif à l'extérieur des
        racines. Pour $n\ge 20$ : $\res{n\ge 41}$.
\end{questions}
\end{exo}

\end{document}
